% !PRINTED X
\weberpage{X}
\inhaltpar{22}{Unterdeterminanten}{72}
\inhaltpar{23}{Die Unterdeterminanten im weiteren Sinne}{77}
\inhaltpar{24}{Lineare homogene Gleichungen}{81}
\inhaltpar{25}{Elimination aus linearen Gleichungen}{87}
\inhaltpar{26}{Unhomogene lineare Gleichungen}{89}
\inhaltpar{27}{Multiplikation von Determinanten}{92}
\inhaltpar{28}{Determinanten der Unterdeterminanten}{95}
\inhaltpar{29}{Interpolation}{98}
\inhaltabschnitt{Dritter Abschnitt.}{Die Wurzeln algebraischer Gleichungen.}
\inhaltpar{30}{Begriff der Wurzeln. Mehrfache Wurzeln}{101}
\inhaltpar{31}{Stetigkeit ganzer Functionen}{103}
\inhaltpar{32}{Vorzeichenwechsel von $f(x)$. Wurzeln von Gleichungen ungeraden Grades und von reinen Gleichungen}{107}
\inhaltpar{33}{Lösung reiner Gleichungen durch trigonometrische Functionen}{111}
\inhaltpar{34}{Befreiung einer Gleichung vom zweiten Gliede}{114}
\inhaltpar{35}{Cubische Gleichungen. \gesperrt{Cardanische} Formel}{116}
\inhaltpar{36}{Der \gesperrt{Cayley}'sche Ausdruck der \gesperrt{Cardanischen} Formel}{118}
\inhaltpar{37}{Die biquadratischen Gleichungen}{119}
\inhaltpar{38}{Beweis des Fundamentalsatzes der Algebra}{121}
\inhaltpar{39}{Algorithmus zur Berechnung der Wurzeln}{127}
\inhaltpar{40}{Stetigkeit der Wurzeln}{132}
\inhaltabschnitt{Vierter Abschnitt.}{Symmetrische Functionen.}
\inhaltpar{41}{Begriff der symmetrischen Functionen. Symmetrische Grundfunctionen}{138}
\inhaltpar{42}{Die Potenzsummen}{140}
\inhaltpar{43}{Beweis des Hauptsatzes für zwei Variable}{143}
\inhaltpar{44}{Allgemeiner Beweis des Hauptsatzes}{144}
\inhaltpar{45}{Zweiter Beweis des Satzes von den symmetrischen Functionen}{147}
\inhaltpar{46}{Discriminanten}{150}
\inhaltpar{47}{Discriminanten der Formen dritter und vierter Ordnung}{153}
\inhaltpar{48}{Resultanten}{156}
\inhaltpar{49}{Elimination. Theorem von \gesperrt{Bezout}}{159}
\inhaltpar{50}{Elimination aus mehreren Gleichungen}{161}
\inhaltpar{51}{Zerlegbare und unzerlegbare Functionen}{164}
\inhaltpar{52}{Tschirnhausen-Transformation}{170}
\inhaltpar{53}{Anwendung auf die cubischen und biquadratischen Gleichungen}{173}
\inhaltpar{54}{Die Tschirnhausen-Transformation der Gleichung 5\textsuperscript{ten} Grades}{175}
\inhaltabschnitt{Fünfter Abschnitt.}{Lineare Transformation. Invarianten.}
\inhaltpar{55}{Lineare Tansformation}{178}
\inhaltpar{56}{Quadratische Formen}{179}
